\documentclass[a4paper]{report}

% begin config
\author{hephaisto}
\newcommand{\moduleTitle}{Gates}
\newcommand{\moduleRevision}{1}
\newcommand{\moduleId}{b}
% end config

\title{\moduleTitle}
\usepackage{graphicx}
\usepackage[a4paper]{geometry}

\usepackage{fancyhdr}
\pagestyle{fancy}
\lhead{BUMM defusal manual}
\rhead{ {\moduleTitle} module (Rev. \moduleRevision)}
\cfoot{}

\renewcommand{\thefigure}{\moduleId.\arabic{figure}} 
\renewcommand{\thetable}{\moduleId.\arabic{table}} 

\usepackage{tikz}
\usepackage{array}

\begin{document}

% begin content

\begin{figure}[t]
	\begin{center}
		\begin{tikzpicture}[scale=0.03,y=-1cm]
			\input{layout_autogen}
		\end{tikzpicture}
	\end{center}
	\caption{Panel layout}
	\label{fig:b_layout}
\end{figure}

\begin{table}
	\centering
	\begin{tabular}{|cm{1cm}m{1cm}|}\hline
%		Signal & LED & Switch \\\hline
		TRUE & \includegraphics{../common/cliparts/led_on} & \includegraphics{../common/cliparts/switch_down} \\
		FALSE & \includegraphics{../common/cliparts/led_off} & \includegraphics{../common/cliparts/switch_up} \\\hline
	\end{tabular}
	\caption{
	\en{Logical states corresponding to input and output elements in default configuration}
	\de{Logischer status der Eingabe- und Ausgabeelemente}
	\chde{Der logische Zustand der Ein- und Ausgabeeinheit}
	}
	\label{tab:b_io_logic}
\end{table}

% BEGIN EXPLANATION 1 (NAND/NOR gates)

\input{manual_figures_autogen}

\section*{\en{How to defuse}\de{Entschärfungsanleitung}\chde{Entschärfungsanleitung}}
\en{
The Gates-module is an accumulation of multiple logical gates connected to other gates and named signal lines.
Each gate has two inputs (binary gate), which can have values of TRUE or FALSE.
Each gate drives all outgoing signal lines to the output value according to the inputs and table \ref{tab:b_truthtable}.
Some named signal lines are connected to LEDs, enabling the user to see the current state of the line.
To defuse this module, set all switches so that they have the value of the corresponding signal line, and press RUN down.

This module is also available with inverted LED values (choose \ledInversionNumbers ) and inverted switch directions (table \ref{tab:b_switch_inversion}).
}
\de{
Das ``Gates'' Modul ist eine Ansammlung von logischen Gattern, die untereinander und mit Signalleitungen verbunden sind.
Jedes Gatter hat zwei Eingabeleitungen, die einen Wert von TRUE (wahr) oder FALSE (falsch) annehmen können, und legt an allen Ausgabeleitungen einen Wert entsprechend Tabelle \ref{tab:b_truthtable} an.
Manche Signalleitungen sind mit LEDs verbunden, sodass man an deren Status den Wert der Leitung erkennen kann.
Um dieses Modul zu entschärfen müssen alle Schalter in die Stellung gebracht werden, die dem Wert der gleichbenannten Signalleitung entspricht, gefolgt von einem Druck auf den RUN Knopf.
Dieses Modul ist auch mit invertierten LED Anzeigen (Seriennummern \ledInversionNumbers ) und invertierten Schalter-Richtungen verfügbar (siehe Tabelle \ref{tab:b_switch_inversion}).
}
\chde{
In ``Gates'' Logikgatter -Modul, mit einander, und sie sind mit der Signalleitung Sammlung verbunden.
Jedes Tor hat zwei Eingangsleitungen , TRUE (true) oder falsch (false) Wert akzeptiert werden kann, und gilt für alle Ausgangsleitungen entsprechend der Wertetabelle \ref{tab:b_truthtable}.
Ein Teil der Signalleitung an die LED-Leuchten angeschlossen, so dass Sie den Wert des Rohres, um ihre Position erkennen kann.
Die Notwendigkeit, alle des Schaltmoduls lösen wird in die Position gebracht, um die Signalleitung gleichbenannten Wert, drücken Sie die Taste RUN entspricht.
Das Modul auch invertierte LED-Anzeige (Seriennummer \ledInversionNumbers) und das invertierte Schaltrichtung (siehe Tabelle \ref{tab:b_switch_inversion}).
}

\begin{table}
	\centering
	\begin{tabular}{|cc|cc|}\hline
		IN 1 & IN 2 & nor & nand \\\hline
		FALSE & FALSE & TRUE & TRUE \\
		FALSE & TRUE & FALSE & TRUE \\
		TRUE & FALSE & FALSE & TRUE \\
		TRUE & TRUE & FALSE & FALSE \\\hline
	\end{tabular}
	\caption{
	\en{Truth table for NAND and NOR gates for different input values}
	\de{Wahrheitstabelle für NAND und NOR Gatter für verschiedene Eingabewerte}
	\chde{Für NAND- und NOR-Gatter für die verschiedenen Eingangswerte Wahrheitstabelle}
	}
	\label{tab:b_truthtable}
\end{table}

% END EXPLANATION 1 (NAND/NOR gates)

\clearpage

% BEGIN EXPLANATION 2 (rules)
\section*{\en{How to defuse}\de{Entschärfungsanleitung}\chde{Entschärfungsanleitung}}
\en{This module can be defused following some simple rules listed below. Be careful, not all rules are equally important and you will have to ignore some! Start with the one rule that is not referenced by others (table \ref{tab:root_rules}). Press the RUN button down when all necessary rules are followed.}
\de{Dieses Modul kann entschärft werden, indem nachfolgende simple Regeln befolgt werden. Nicht alle Regeln sind wichtig, viele müssen ignoriert werden! Starte mit der einen Regel, die von keiner anderen erwähnt wird (siehe Tabelle \ref{tab:root_rules}). Drücke den RUN-Knopf herunter, wenn alle relevanten Regeln erfüllt sind.}
\chde{Das Modul kann durch die folgenden einfachen Regeln erhalten werden, zur Einhaltung erleichtert werden. Nicht alle Regeln sind wichtig, viele müssen ignoriert werden! Beginnen Sie mit einem anderen Referenz (siehe Tabelle \ref{tab:root_rules}) Regeln. Wenn Sie die RUN-Taste drücken, werden alle einschlägigen Vorschriften erfüllt sind.}

\en{This module is also available with inverted LED values (choose \ledInversionNumbers ) and inverted switch directions (table \ref{tab:b_switch_inversion}).}
\de{Dieses Modul ist auch mit invertierten LED Anzeigen (Seriennummern \ledInversionNumbers ) und invertierten Schalter-Richtungen verfügbar (siehe Tabelle \ref{tab:b_switch_inversion}).}
\chde{Das Modul auch invertierte LED-Anzeige (Seriennummer \ledInversionNumbers) und das invertierte Schaltrichtung (siehe Tabelle \ref{tab:b_switch_inversion}).}
\input{flowchoice_autogen}
% end EXPLANATION 2 (rules)

% end content

\end{document}
